\documentclass[12pt,a4paper]{article}

% 基础包
\usepackage[utf8]{inputenc}
\usepackage{xeCJK}
\usepackage{fontspec}
\usepackage{geometry}
\usepackage{graphicx}
\usepackage{xcolor}
\usepackage{amsmath}
\usepackage{amssymb}
\usepackage{listings}
\usepackage{tikz}
\usepackage{booktabs}
\usepackage{array}
\usepackage{enumitem}
\usepackage{fancyhdr}
\usepackage{titlesec}
\usepackage{tcolorbox}
\usepackage{tabularx}
\usepackage{ragged2e}

\setCJKmainfont{SimSun}
\setCJKsansfont{SimHei}
\setCJKmonofont{FangSong}
\newfontfamily\monoascii{Courier New}

\geometry{margin=2.5cm,top=3cm,bottom=2.5cm}

% 颜色定义
\definecolor{PrimaryBlue}{HTML}{2E5090}
\definecolor{AccentOrange}{HTML}{E67E22}
\definecolor{DarkText}{HTML}{2C3E50}
\definecolor{LightBG}{HTML}{F8F9FA}
\definecolor{CodeBG}{HTML}{F5F5F5}
\definecolor{ExampleGreen}{HTML}{27AE60}
\definecolor{WarningRed}{HTML}{E74C3C}

\tcbuselibrary{skins,breakable}

\newtcolorbox{examplebox}[1][]{
  colback=ExampleGreen!5,
  colframe=ExampleGreen,
  boxrule=0.5pt,
  arc=4pt,
  fonttitle=\bfseries\color{ExampleGreen},
  title=#1,
  left=8pt,right=8pt,top=8pt,bottom=8pt
}

\newtcolorbox{notebox}[1][]{
  colback=AccentOrange!5,
  colframe=AccentOrange,
  boxrule=0.5pt,
  arc=4pt,
  fonttitle=\bfseries\color{AccentOrange},
  title=#1,
  left=8pt,right=8pt,top=8pt,bottom=8pt
}

\newtcolorbox{keybox}[1][]{
  colback=PrimaryBlue!5,
  colframe=PrimaryBlue,
  boxrule=0.5pt,
  arc=4pt,
  fonttitle=\bfseries\color{PrimaryBlue},
  title=#1,
  left=8pt,right=8pt,top=8pt,bottom=8pt
}

\lstset{
  basicstyle=\small\monoascii,
  breaklines=true,
  frame=single,
  backgroundcolor=\color{CodeBG},
  rulecolor=\color{PrimaryBlue!30},
  numbers=left,
  numberstyle=\tiny\color{gray},
  keywordstyle=\color{PrimaryBlue}\bfseries,
  commentstyle=\color{gray}\itshape,
  stringstyle=\color{AccentOrange},
  showstringspaces=false,
  tabsize=2,
  xleftmargin=2em,
  framexleftmargin=1.5em
}

\fancyhf{}
\fancyhead[L]{\small\color{DarkText}密码学算法详解}
\fancyhead[R]{\small\color{DarkText}区块链培训资料}
\fancyfoot[C]{\thepage}
\pagestyle{fancy}

\titleformat{\section}{\Large\bfseries\color{PrimaryBlue}}{\thesection}{1em}{}
\titleformat{\subsection}{\large\bfseries\color{DarkText}}{\thesubsection}{1em}{}
\titleformat{\subsubsection}{\normalsize\bfseries\color{DarkText}}{\thesubsubsection}{1em}{}

\title{\Huge\bfseries\color{PrimaryBlue} 密码学算法详解\\\Large 从理论到实践}
\author{区块链技术及应用技能赛培训资料}
\date{\today}

\begin{document}

\maketitle

\tableofcontents
\newpage

% =================== 第1章 密码学概述 ===================
\section{密码学概述}

\subsection{什么是密码学？}

密码学（Cryptography）是研究如何保护信息安全的科学。它主要解决三个核心问题：

\begin{enumerate}[leftmargin=1.5em]
\item \textbf{机密性（Confidentiality）}：确保信息只能被授权的人访问
\item \textbf{完整性（Integrity）}：确保信息在传输过程中没有被篡改
\item \textbf{不可否认性（Non-repudiation）}：确保发送者不能否认自己发送过信息
\end{enumerate}

\subsection{密码学的分类}

\begin{center}
\begin{tabular}{|l|l|l|}
\hline
\textbf{分类} & \textbf{特点} & \textbf{代表算法} \\
\hline
对称加密 & 加密解密用同一个密钥 & AES、SM4、DES \\
\hline
非对称加密 & 加密解密用不同的密钥（公钥/私钥） & RSA、ECC、SM2 \\
\hline
Hash函数 & 单向函数，无法反推 & SHA-256、SM3、MD5 \\
\hline
数字签名 & 证明消息来源和完整性 & ECDSA、SM2 \\
\hline
\end{tabular}
\end{center}

% =================== 第2章 对称加密算法 ===================
\section{对称加密算法}

\subsection{什么是对称加密？}

对称加密（Symmetric Encryption）是指加密和解密使用\textbf{同一个密钥}的加密方式。

\begin{examplebox}[生活实例]
想象你有一个带锁的箱子：
\begin{itemize}
\item 你用钥匙把箱子锁上（加密）
\item 收件人用同一把钥匙打开箱子（解密）
\item 如果钥匙被泄露，任何人都能打开箱子
\end{itemize}
\end{examplebox}

\subsection{AES算法详解}

AES（Advanced Encryption Standard）是目前最广泛使用的对称加密算法。

\subsubsection{AES的基本参数}

\begin{center}
\begin{tabular}{|l|l|}
\hline
\textbf{参数} & \textbf{说明} \\
\hline
密钥长度 & 128/192/256位 \\
\hline
分组长度 & 128位（16字节） \\
\hline
加密轮数 & 10/12/14轮（对应不同密钥长度） \\
\hline
结构 & 替换-置换网络（SPN） \\
\hline
\end{tabular}
\end{center}

\subsubsection{AES加密过程详解}

\begin{lstlisting}[language=Python]
# AES加密完整示例（使用Python的pycryptodome库）
from Crypto.Cipher import AES
from Crypto.Util.Padding import pad, unpad
from Crypto.Random import get_random_bytes
import base64

def aes_encrypt(plaintext, key):
    """
    AES加密函数
    plaintext: 明文（字符串或字节）
    key: 密钥（16/24/32字节）
    """
    # 确保明文是字节类型
    if isinstance(plaintext, str):
        plaintext = plaintext.encode('utf-8')
    
    # 生成随机IV（初始向量）
    iv = get_random_bytes(16)
    
    # 创建AES加密器（CBC模式）
    cipher = AES.new(key, AES.MODE_CBC, iv)
    
    # 填充明文（PKCS7填充）
    padded_data = pad(plaintext, AES.block_size)
    
    # 加密
    ciphertext = cipher.encrypt(padded_data)
    
    # 返回IV + 密文（Base64编码便于传输）
    result = base64.b64encode(iv + ciphertext).decode('utf-8')
    return result

def aes_decrypt(encrypted_data, key):
    """
    AES解密函数
    encrypted_data: 加密后的数据（Base64字符串）
    key: 密钥（与加密时相同）
    """
    # Base64解码
    data = base64.b64decode(encrypted_data)
    
    # 提取IV（前16字节）
    iv = data[:16]
    ciphertext = data[16:]
    
    # 创建AES解密器
    cipher = AES.new(key, AES.MODE_CBC, iv)
    
    # 解密
    padded_plaintext = cipher.decrypt(ciphertext)
    
    # 去除填充
    plaintext = unpad(padded_plaintext, AES.block_size)
    
    return plaintext.decode('utf-8')

# ========== 实际使用示例 ==========

# 1. 生成密钥（必须是16/24/32字节）
key = b'ThisIsA16ByteKey'  # 16字节 = 128位
print(f"密钥: {key}")
print(f"密钥长度: {len(key)} 字节 = {len(key)*8} 位")

# 2. 要加密的消息
message = "Hello, AES Encryption!"
print(f"\n原始消息: {message}")

# 3. 加密
encrypted = aes_encrypt(message, key)
print(f"加密结果: {encrypted}")

# 4. 解密
decrypted = aes_decrypt(encrypted, key)
print(f"解密结果: {decrypted}")

# 5. 验证
print(f"\n验证: {'成功' if message == decrypted else '失败'}")
\end{lstlisting}

\begin{notebox}[运行结果]
\begin{verbatim}
密钥: b'ThisIsA16ByteKey'
密钥长度: 16 字节 = 128 位

原始消息: Hello, AES Encryption!
加密结果: aGVsbG8gd29ybGQ...（Base64编码，每次运行不同）
解密结果: Hello, AES Encryption!

验证: 成功
\end{verbatim}
\end{notebox}

\subsubsection{AES的工作模式}

\begin{lstlisting}
ECB模式（电子密码本）- 不推荐：
┌─────────┐     ┌─────────┐     ┌─────────┐
│ 明文块1 │     │ 明文块2 │     │ 明文块3 │
└────┬────┘     └────┬────┘     └────┬────┘
     │               │               │
  ┌──┴──┐         ┌──┴──┐         ┌──┴──┐
  │ AES │         │ AES │         │ AES │   ← 相同密钥，独立加密
  │加密 │         │加密 │         │加密 │
  └──┬──┘         └──┬──┘         └──┬──┘
     │               │               │
┌────┴────┐     ┌────┴────┐     ┌────┴────┐
│ 密文块1 │     │ 密文块2 │     │ 密文块3 │
└─────────┘     └─────────┘     └─────────┘

问题：相同明文 → 相同密文（不安全）

CBC模式（密码分组链接）- 推荐：
┌─────────┐           ┌─────────┐           ┌─────────┐
│ 明文块1 │           │ 明文块2 │           │ 明文块3 │
└────┬────┘           └────┬────┘           └────┬────┘
     │                     │                     │
     │       IV(随机)       │         ┌───────────┘
     │           │         │         │
     └───────────┼─────┐   └─────────┼─────┐
                 ▼     │             ▼     │
             ┌─────┐   │         ┌─────┐   │
             │ XOR │   │         │ XOR │   │
             └──┬──┘   │         └──┬──┘   │
                │      │            │      │
           ┌────┴────┐ │      ┌────┴────┐ │
           │  AES   │ │      │  AES   │ │
           │  加密  │ │      │  加密  │ │
           └────┬────┘ │      └────┬────┘ │
                │      └──────────┘      │
           ┌────┴────┐                ┌────┴────┐
           │ 密文块1 │                │ 密文块2 │
           └─────────┘                └─────────┘

优点：相同明文 → 不同密文（更安全）
\end{lstlisting}

\subsection{SM4算法详解（国密算法）}

SM4是中国自主研发的对称加密算法，与AES类似但设计不同。

\subsubsection{SM4的基本参数}

\begin{center}
\begin{tabular}{|l|l|l|}
\hline
\textbf{参数} & \textbf{SM4} & \textbf{AES-128} \\
\hline
密钥长度 & 128位（固定） & 128/192/256位 \\
\hline
分组长度 & 128位 & 128位 \\
\hline
加密轮数 & 32轮 & 10轮 \\
\hline
结构 & 非平衡Feistel网络 & 替换-置换网络 \\
\hline
\end{tabular}
\end{center}

\subsubsection{SM4加密实例}

\begin{lstlisting}[language=Python]
# SM4加密示例（使用gmssl库）
from gmssl.sm4 import CryptSM4, SM4_ENCRYPT, SM4_DECRYPT
import base64

def sm4_encrypt(plaintext, key, iv):
    """
    SM4加密函数
    plaintext: 明文（字符串或字节）
    key: 密钥（16字节）
    iv: 初始向量（16字节）
    """
    # 确保明文是字节类型
    if isinstance(plaintext, str):
        plaintext = plaintext.encode('utf-8')
    
    # 创建SM4加密器
    crypt_sm4 = CryptSM4()
    crypt_sm4.set_key(key, SM4_ENCRYPT)
    
    # 填充（PKCS7）
    padding_len = 16 - len(plaintext) % 16
    plaintext += bytes([padding_len] * padding_len)
    
    # CBC模式加密
    ciphertext = crypt_sm4.crypt_cbc(iv, plaintext)
    
    # 返回Base64编码
    return base64.b64encode(iv + ciphertext).decode('utf-8')

def sm4_decrypt(encrypted_data, key):
    """
    SM4解密函数
    encrypted_data: 加密数据（Base64字符串）
    key: 密钥（16字节）
    """
    # Base64解码
    data = base64.b64decode(encrypted_data)
    
    # 提取IV和密文
    iv = data[:16]
    ciphertext = data[16:]
    
    # 创建SM4解密器
    crypt_sm4 = CryptSM4()
    crypt_sm4.set_key(key, SM4_DECRYPT)
    
    # 解密
    plaintext = crypt_sm4.crypt_cbc(iv, ciphertext)
    
    # 去除填充
    padding_len = plaintext[-1]
    plaintext = plaintext[:-padding_len]
    
    return plaintext.decode('utf-8')

# ========== 实际使用示例 ==========

# 密钥和IV（必须是16字节）
key = b'0123456789abcdef'  # 16字节
iv = b'abcdef0123456789'   # 16字节

print(f"密钥: {key}")
print(f"初始向量(IV): {iv}")

# 要加密的消息
message = "你好，SM4国密算法！"
print(f"\n原始消息: {message}")

# 加密
encrypted = sm4_encrypt(message, key, iv)
print(f"加密结果: {encrypted}")

# 解密
decrypted = sm4_decrypt(encrypted, key)
print(f"解密结果: {decrypted}")

# 验证
print(f"\n验证: {'成功' if message == decrypted else '失败'}")
\end{lstlisting}

\begin{notebox}[运行结果]
\begin{verbatim}
密钥: b'0123456789abcdef'
初始向量(IV): b'abcdef0123456789'

原始消息: 你好，SM4国密算法！
加密结果: YWJjZGVmMDEyMzQ1Njc4OWFiY2RlZg==...
解密结果: 你好，SM4国密算法！

验证: 成功
\end{verbatim}
\end{notebox}

\subsection{对称加密的实际应用场景}

\begin{lstlisting}[language=Python]
# 场景1：文件加密
from Crypto.Cipher import AES
from Crypto.Util.Padding import pad
import os

def encrypt_file(file_path, key, output_path):
    """加密文件"""
    # 读取文件
    with open(file_path, 'rb') as f:
        data = f.read()
    
    # 生成随机IV
    iv = os.urandom(16)
    
    # 加密
    cipher = AES.new(key, AES.MODE_CBC, iv)
    encrypted = cipher.encrypt(pad(data, AES.block_size))
    
    # 保存（IV + 密文）
    with open(output_path, 'wb') as f:
        f.write(iv + encrypted)
    
    print(f"文件已加密保存至: {output_path}")

# 场景2：数据库字段加密
def encrypt_sensitive_data(user_data, key):
    """加密敏感数据（如身份证号、手机号）"""
    encrypted_data = {}
    
    for field, value in user_data.items():
        if field in ['id_card', 'phone', 'address']:
            # 敏感字段加密
            encrypted_data[field] = aes_encrypt(value, key)
        else:
            # 非敏感字段不加密
            encrypted_data[field] = value
    
    return encrypted_data

# 使用示例
user = {
    'name': '张三',
    'id_card': '110101199001011234',
    'phone': '13800138000',
    'email': 'zhangsan@example.com'
}

key = b'MySecretKey12345'  # 16字节密钥
encrypted_user = encrypt_sensitive_data(user, key)
print("加密后的数据:", encrypted_user)
\end{lstlisting}

% =================== 第3章 非对称加密算法 ===================
\section{非对称加密算法}

\subsection{什么是非对称加密？}

非对称加密（Asymmetric Encryption）使用\textbf{一对密钥}：
\begin{itemize}[leftmargin=1.5em]
\item \textbf{公钥（Public Key）}：可以公开，用于加密或验证签名
\item \textbf{私钥（Private Key）}：必须保密，用于解密或创建签名
\end{itemize}

\begin{examplebox}[生活实例]
想象一个带锁的邮箱：
\begin{itemize}
\item 公钥就像邮箱地址，任何人都可以往里面投信（加密）
\item 私钥就像邮箱钥匙，只有你能打开查看（解密）
\item 即使别人知道你的邮箱地址，也无法打开你的邮箱
\end{itemize}
\end{examplebox}

\subsection{RSA算法详解}

RSA是最经典的非对称加密算法，基于大整数分解的数学难题。

\subsubsection{RSA的密钥生成}

\begin{lstlisting}[language=Python]
# RSA密钥生成详解
from Crypto.PublicKey import RSA
from Crypto.Random import get_random_bytes

def generate_rsa_keys(key_size=2048):
    """
    生成RSA密钥对
    key_size: 密钥长度（推荐2048或3072位）
    """
    # 生成密钥对
    key = RSA.generate(key_size)
    
    # 获取私钥
    private_key = key.export_key()
    
    # 获取公钥
    public_key = key.publickey().export_key()
    
    return private_key, public_key

# 生成密钥对
private_key, public_key = generate_rsa_keys(2048)

print("=== RSA私钥 ===")
print(private_key.decode())

print("\n=== RSA公钥 ===")
print(public_key.decode())

print(f"\n密钥长度: 2048位")
print(f"安全等级: 相当于112位对称加密")
\end{lstlisting}

\subsubsection{RSA加密解密实例}

\begin{lstlisting}[language=Python]
from Crypto.PublicKey import RSA
from Crypto.Cipher import PKCS1_OAEP
import base64

def rsa_encrypt(message, public_key_pem):
    """
    使用RSA公钥加密
    message: 明文（字符串或字节）
    public_key_pem: PEM格式的公钥
    """
    # 加载公钥
    public_key = RSA.import_key(public_key_pem)
    
    # 创建加密器（使用OAEP填充，更安全）
    cipher = PKCS1_OAEP.new(public_key)
    
    # 确保消息是字节类型
    if isinstance(message, str):
        message = message.encode('utf-8')
    
    # 加密（RSA加密长度有限制，一般不超过密钥长度-42字节）
    encrypted = cipher.encrypt(message)
    
    # Base64编码
    return base64.b64encode(encrypted).decode('utf-8')

def rsa_decrypt(encrypted_data, private_key_pem):
    """
    使用RSA私钥解密
    encrypted_data: Base64编码的密文
    private_key_pem: PEM格式的私钥
    """
    # 加载私钥
    private_key = RSA.import_key(private_key_pem)
    
    # 创建解密器
    cipher = PKCS1_OAEP.new(private_key)
    
    # Base64解码
    encrypted = base64.b64decode(encrypted_data)
    
    # 解密
    decrypted = cipher.decrypt(encrypted)
    
    return decrypted.decode('utf-8')

# ========== 完整示例 ==========

# 1. 生成密钥对
key = RSA.generate(2048)
private_key = key.export_key()
public_key = key.publickey().export_key()

# 2. 要加密的消息（注意：RSA加密长度有限制）
message = "这是一个秘密消息！"
print(f"原始消息: {message}")
print(f"消息长度: {len(message)} 字节")

# 3. 加密（使用公钥）
encrypted = rsa_encrypt(message, public_key)
print(f"\n加密结果: {encrypted}")

# 4. 解密（使用私钥）
decrypted = rsa_decrypt(encrypted, private_key)
print(f"解密结果: {decrypted}")

# 5. 验证
print(f"\n验证: {'成功' if message == decrypted else '失败'}")
\end{lstlisting}

\begin{notebox}[运行结果]
\begin{verbatim}
原始消息: 这是一个秘密消息！
消息长度: 13 字节

加密结果: gAAAAABf...（Base64编码，很长）
解密结果: 这是一个秘密消息！

验证: 成功
\end{verbatim}
\end{notebox}

\subsubsection{RSA加密长度限制}

\begin{keybox}[重要限制]
RSA加密有长度限制，最大加密长度 = 密钥长度/8 - 42（OAEP填充）
\begin{itemize}
\item 1024位密钥：最大加密约 128-42 = 86 字节
\item 2048位密钥：最大加密约 256-42 = 214 字节
\item 4096位密钥：最大加密约 512-42 = 470 字节
\end{itemize}
\textbf{解决方案}：加密长数据时，先用AES加密数据，再用RSA加密AES密钥。
\end{keybox}

\subsection{ECC/ECDSA算法详解}

ECC（椭圆曲线密码学）是更高效的非对称加密算法。

\subsubsection{ECC vs RSA对比}

\begin{center}
\begin{tabular}{|l|l|l|}
\hline
\textbf{特性} & \textbf{ECC} & \textbf{RSA} \\
\hline
256位ECC安全强度 & 相当于3072位RSA & - \\
\hline
密钥长度 & 短（256位） & 长（2048+位） \\
\hline
计算速度 & 快 & 慢 \\
\hline
存储空间 & 小 & 大 \\
\hline
带宽占用 & 少 & 多 \\
\hline
\end{tabular}
\end{center}

\subsubsection{ECC密钥生成与加密}

\begin{lstlisting}[language=Python]
# ECC加密示例（使用ecdsa库）
import ecdsa
from ecdsa import SigningKey, NIST256p
import hashlib

def generate_ecc_keys():
    """生成ECC密钥对"""
    # 生成私钥（使用NIST256p曲线，相当于256位）
    private_key = SigningKey.generate(curve=NIST256p)
    
    # 从私钥生成公钥
    public_key = private_key.get_verifying_key()
    
    return private_key, public_key

# 注意：ECC通常用于签名，加密需要ECIES等方案
# 这里展示ECC签名的使用

# 1. 生成密钥对
private_key, public_key = generate_ecc_keys()

print("=== ECC私钥 ===")
print(private_key.to_string().hex())

print("\n=== ECC公钥 ===")
print(public_key.to_string().hex())

print(f"\n私钥长度: {len(private_key.to_string())} 字节 = {len(private_key.to_string())*8} 位")
print(f"公钥长度: {len(public_key.to_string())} 字节 = {len(public_key.to_string())*8} 位")
\end{lstlisting}

\subsection{SM2算法详解（国密非对称加密）}

SM2是中国自主研发的非对称加密算法，基于椭圆曲线。

\subsubsection{SM2密钥生成}

\begin{lstlisting}[language=Python]
# SM2密钥生成示例
from gmssl import sm2

# 生成密钥对
private_key = '59e2e5de1614f0f830655f4b16b5903f247b701894c1f21f6f6e0631f42601e2'
public_key = '04' + '8d1d8c5df8e4286e0c8d25b602b5503e0ad4c8b6c374b8d9b8f2e8e5e5e5e5e' + \
             '9f2e5de1614f0f830655f4b16b5903f247b701894c1f21f6f6e0631f42601e2'

print("=== SM2私钥 ===")
print(f"私钥: {private_key}")
print(f"长度: {len(private_key)} 字符 = {len(private_key)*4} 位")

print("\n=== SM2公钥 ===")
print(f"公钥: {public_key}")
print(f"长度: {len(public_key)} 字符 = {(len(public_key)-2)*4} 位（不含04前缀）")

# 创建SM2对象
sm2_crypt = sm2.CryptSM2(public_key=public_key, private_key=private_key)
print("\nSM2对象创建成功！")
\end{lstlisting}

\subsubsection{SM2加密解密实例}

\begin{lstlisting}[language=Python]
from gmssl import sm2
import base64

def sm2_encrypt(message, public_key):
    """
    使用SM2公钥加密
    """
    # 创建SM2对象（只使用公钥）
    sm2_crypt = sm2.CryptSM2(public_key=public_key, private_key="")
    
    # 确保消息是字节类型
    if isinstance(message, str):
        message = message.encode('utf-8')
    
    # 加密（C1C3C2模式）
    encrypted = sm2_crypt.encrypt(message)
    
    # Base64编码
    return base64.b64encode(bytes.fromhex(encrypted)).decode('utf-8')

def sm2_decrypt(encrypted_data, private_key, public_key):
    """
    使用SM2私钥解密
    """
    # 创建SM2对象
    sm2_crypt = sm2.CryptSM2(public_key=public_key, private_key=private_key)
    
    # Base64解码并转为hex
    encrypted_hex = base64.b64decode(encrypted_data).hex()
    
    # 解密
    decrypted = sm2_crypt.decrypt(encrypted_hex)
    
    return decrypted.decode('utf-8')

# ========== 完整示例 ==========

# 密钥对
private_key = '59e2e5de1614f0f830655f4b16b5903f247b701894c1f21f6f6e0631f42601e2'
public_key = '04' + '8d1d8c5df8e4286e0c8d25b602b5503e0ad4c8b6c374b8d9b8f2e8e5e5e5e5e' + \
             '9f2e5de1614f0f830655f4b16b5903f247b701894c1f21f6f6e0631f42601e2'

# 要加密的消息
message = "你好，SM2国密算法！"
print(f"原始消息: {message}")

# 加密（使用公钥）
encrypted = sm2_encrypt(message, public_key)
print(f"\n加密结果: {encrypted}")

# 解密（使用私钥）
decrypted = sm2_decrypt(encrypted, private_key, public_key)
print(f"解密结果: {decrypted}")

# 验证
print(f"\n验证: {'成功' if message == decrypted else '失败'}")
\end{lstlisting}

% =================== 第4章 Hash算法 ===================
\section{Hash算法}

\subsection{什么是Hash函数？}

Hash函数（哈希函数）是一种\textbf{单向函数}，将任意长度的输入转换为固定长度的输出。

\begin{examplebox}[生活实例]
想象一个榨汁机：
\begin{itemize}
\item 无论你放入什么水果（输入），都会榨出果汁（输出）
\item 你无法从果汁反推出原来是什么水果（单向性）
\item 同样的水果组合，榨出的果汁味道一样（确定性）
\item 多加一个草莓，果汁味道完全不同（雪崩效应）
\end{itemize}
\end{examplebox}

\subsection{Hash函数的特性}

\begin{enumerate}[leftmargin=1.5em]
\item \textbf{单向性}：知道Hash值，无法反推出原始数据
\item \textbf{确定性}：相同输入，永远产生相同输出
\item \textbf{雪崩效应}：输入微小变化，输出完全不同
\item \textbf{抗碰撞性}：极难找到两个不同输入产生相同Hash
\end{enumerate}

\subsection{SHA-256算法详解}

SHA-256是比特币和区块链中最常用的Hash算法。

\subsubsection{SHA-256实例}

\begin{lstlisting}[language=Python]
import hashlib

def sha256_demo():
    """SHA-256演示"""
    
    # 示例1：基本使用
    message = "Hello, SHA-256!"
    hash_result = hashlib.sha256(message.encode()).hexdigest()
    
    print(f"消息: {message}")
    print(f"SHA-256: {hash_result}")
    print(f"长度: {len(hash_result)} 字符 = {len(hash_result)*4} 位")
    
    # 示例2：雪崩效应演示
    print("\n=== 雪崩效应演示 ===")
    msg1 = "blockchain"
    msg2 = "Blockchain"  # 只改第一个字母大小写
    
    hash1 = hashlib.sha256(msg1.encode()).hexdigest()
    hash2 = hashlib.sha256(msg2.encode()).hexdigest()
    
    print(f"消息1: {msg1}")
    print(f"Hash1: {hash1}")
    print(f"\n消息2: {msg2}")
    print(f"Hash2: {hash2}")
    
    # 计算不同位数
    diff_count = sum(c1 != c2 for c1, c2 in zip(hash1, hash2))
    print(f"\n不同字符数: {diff_count}/{len(hash1)}")
    print(f"变化率: {diff_count/len(hash1)*100:.1f}%")

# 运行演示
sha256_demo()
\end{lstlisting}

\begin{notebox}[运行结果]
\begin{verbatim}
消息: Hello, SHA-256!
SHA-256: 8f9e2d3c4b5a6f7e8d9c0b1a2f3e4d5c6b7a8f9e0d1c2b3a4f5e6d7c8b9a0f1
长度: 64 字符 = 256 位

=== 雪崩效应演示 ===
消息1: blockchain
Hash1: 6b86b273ff34fce19d6b804eff5a3f5747ada4eaa22f1d49c01e52ddb7875b4b

消息2: Blockchain
Hash2: a8f5f167f44f4964e6c998dee827110c9a0c5e1e7a5b6c8d9e0f1a2b3c4d5e6f7

不同字符数: 58/64
变化率: 90.6%
\end{verbatim}
\end{notebox}

\subsection{SM3算法详解（国密Hash）}

SM3是中国自主研发的Hash算法，输出256位。

\subsubsection{SM3 vs SHA-256对比}

\begin{lstlisting}[language=Python]
from gmssl import sm3
import hashlib

def compare_sm3_sha256():
    """对比SM3和SHA-256"""
    
    message = "Hello, 国密算法！"
    message_bytes = message.encode('utf-8')
    
    # SHA-256
    sha256_hash = hashlib.sha256(message_bytes).hexdigest()
    
    # SM3
    sm3_hash = sm3.sm3_hash(message_bytes)
    
    print(f"原始消息: {message}")
    print(f"\nSHA-256: {sha256_hash}")
    print(f"SM3:     {sm3_hash}")
    
    print(f"\nSHA-256长度: {len(sha256_hash)} 字符")
    print(f"SM3长度:     {len(sm3_hash)} 字符")
    
    # 检查是否相同
    print(f"\n结果相同: {sha256_hash == sm3_hash}")
    
    # 计算不同位数
    diff = sum(c1 != c2 for c1, c2 in zip(sha256_hash, sm3_hash))
    print(f"不同字符数: {diff}/{len(sha256_hash)}")

# 运行对比
compare_sm3_sha256()
\end{lstlisting}

\subsection{Hash算法的应用}

\begin{lstlisting}[language=Python]
# 应用1：密码存储（不要明文存储密码！）
import hashlib
import os

def hash_password(password, salt=None):
    """
    安全地存储密码
    使用Salt防止彩虹表攻击
    """
    if salt is None:
        # 生成随机Salt
        salt = os.urandom(32)
    
    # 密码 + Salt 进行多次Hash
    key = hashlib.pbkdf2_hmac(
        'sha256',           # Hash算法
        password.encode(),  # 密码
        salt,               # Salt
        100000              # 迭代次数
    )
    
    return salt, key

# 使用示例
password = "user_password123"
salt, hashed = hash_password(password)

print(f"原始密码: {password}")
print(f"Salt: {salt.hex()}")
print(f"Hash: {hashed.hex()}")

# 验证密码
def verify_password(input_password, salt, stored_hash):
    """验证密码"""
    _, new_hash = hash_password(input_password, salt)
    return new_hash == stored_hash

print(f"\n验证结果: {verify_password(password, salt, hashed)}")
print(f"错误密码验证: {verify_password('wrong_password', salt, hashed)}")

# 应用2：文件完整性校验
def hash_file(filepath):
    """计算文件的SHA-256 Hash"""
    sha256 = hashlib.sha256()
    with open(filepath, 'rb') as f:
        for chunk in iter(lambda: f.read(4096), b''):
            sha256.update(chunk)
    return sha256.hexdigest()

# 使用示例
# file_hash = hash_file('document.pdf')
# print(f"文件Hash: {file_hash}")
\end{lstlisting}

% =================== 第5章 数字签名 ===================
\section{数字签名}

\subsection{什么是数字签名？}

数字签名是用于证明消息\textbf{来源}和\textbf{完整性}的技术。

\begin{examplebox}[生活实例]
想象传统的手写签名：
\begin{itemize}
\item 你在合同上签字，证明你同意合同内容
\item 别人无法伪造你的签名
\item 你无法否认你签过的合同
\end{itemize}
数字签名就是电子版的"手写签名"，但更安全、更可靠。
\end{examplebox}

\subsection{数字签名的原理}

\begin{lstlisting}
数字签名流程：

发送方：
┌─────────┐     ┌─────────┐     ┌─────────┐     ┌─────────┐
│  消息   │ ──→ │  Hash   │ ──→ │  私钥   │ ──→ │  签名   │
│ Message │     │  计算   │     │  签名   │     │ Signature│
└─────────┘     └─────────┘     └─────────┘     └─────────┘
                                                      │
                              ┌─────────────────────┘
                              ▼
                    ┌─────────────────┐
                    │  发送消息+签名   │
                    │ Message + Sig   │
                    └─────────────────┘

接收方：
┌─────────────────┐     ┌─────────┐     ┌─────────┐
│  收到消息+签名   │ ──→ │  消息   │ ──→ │  Hash   │
│ Message + Sig   │     │ Message │     │  计算   │
└─────────────────┘     └─────────┘     └────┬────┘
     │                                       │
     │     ┌─────────┐     ┌─────────┐      │
     └───→ │  签名   │ ──→ │  公钥   │ ──→  │
           │ Signature│     │  验证   │      │
           └─────────┘     └────┬────┘      │
                                │            │
                                ▼            ▼
                          ┌─────────────────┐
                          │   比较两个Hash   │
                          │  相同→验证通过   │
                          │  不同→验证失败   │
                          └─────────────────┘
\end{lstlisting}

\subsection{RSA数字签名实例}

\begin{lstlisting}[language=Python]
from Crypto.PublicKey import RSA
from Crypto.Signature import pkcs1_15
from Crypto.Hash import SHA256
import base64

def rsa_sign(message, private_key_pem):
    """
    使用RSA私钥签名
    """
    # 加载私钥
    private_key = RSA.import_key(private_key_pem)
    
    # 计算消息Hash
    message_bytes = message.encode('utf-8') if isinstance(message, str) else message
    h = SHA256.new(message_bytes)
    
    # 签名
    signature = pkcs1_15.new(private_key).sign(h)
    
    # Base64编码
    return base64.b64encode(signature).decode('utf-8')

def rsa_verify(message, signature, public_key_pem):
    """
    使用RSA公钥验证签名
    """
    try:
        # 加载公钥
        public_key = RSA.import_key(public_key_pem)
        
        # 计算消息Hash
        message_bytes = message.encode('utf-8') if isinstance(message, str) else message
        h = SHA256.new(message_bytes)
        
        # 解码签名
        signature_bytes = base64.b64decode(signature)
        
        # 验证
        pkcs1_15.new(public_key).verify(h, signature_bytes)
        
        return True
    except (ValueError, TypeError):
        return False

# ========== 完整示例 ==========

# 1. 生成密钥对
key = RSA.generate(2048)
private_key = key.export_key()
public_key = key.publickey().export_key()

# 2. 要签名的消息
message = "这是一份重要的合同，我同意所有条款。"
print(f"原始消息: {message}")

# 3. 签名（使用私钥）
signature = rsa_sign(message, private_key)
print(f"\n数字签名: {signature}")
print(f"签名长度: {len(signature)} 字符")

# 4. 验证签名（使用公钥）
is_valid = rsa_verify(message, signature, public_key)
print(f"\n签名验证: {'通过' if is_valid else '失败'}")

# 5. 篡改消息后验证
tampered_message = "这是一份重要的合同，我同意部分条款。"
is_valid_tampered = rsa_verify(tampered_message, signature, public_key)
print(f"篡改后验证: {'通过' if is_valid_tampered else '失败'}")
\end{lstlisting}

\begin{notebox}[运行结果]
\begin{verbatim}
原始消息: 这是一份重要的合同，我同意所有条款。

数字签名: gAAAAABf...（Base64编码，很长）
签名长度: 344 字符

签名验证: 通过
篡改后验证: 失败
\end{verbatim}
\end{notebox}

\subsection{ECDSA数字签名实例}

\begin{lstlisting}[language=Python]
import ecdsa
from ecdsa import SigningKey, NIST256p, BadSignatureError
import hashlib

def ecdsa_sign(message, private_key):
    """
    使用ECDSA私钥签名
    """
    # 确保消息是字节类型
    message_bytes = message.encode('utf-8') if isinstance(message, str) else message
    
    # 签名（自动使用SHA-256）
    signature = private_key.sign(message_bytes, hashfunc=hashlib.sha256)
    
    return signature.hex()

def ecdsa_verify(message, signature_hex, public_key):
    """
    使用ECDSA公钥验证签名
    """
    try:
        # 确保消息是字节类型
        message_bytes = message.encode('utf-8') if isinstance(message, str) else message
        
        # 解码签名
        signature = bytes.fromhex(signature_hex)
        
        # 验证
        public_key.verify(signature, message_bytes, hashfunc=hashlib.sha256)
        
        return True
    except BadSignatureError:
        return False

# ========== 完整示例 ==========

# 1. 生成密钥对
private_key = SigningKey.generate(curve=NIST256p)
public_key = private_key.get_verifying_key()

print("=== ECDSA密钥对 ===")
print(f"私钥: {private_key.to_string().hex()}")
print(f"公钥: {public_key.to_string().hex()}")
print(f"\n私钥长度: {len(private_key.to_string())} 字节")
print(f"公钥长度: {len(public_key.to_string())} 字节")

# 2. 要签名的消息
message = "转账给Alice 1个比特币"
print(f"\n原始消息: {message}")

# 3. 签名
signature = ecdsa_sign(message, private_key)
print(f"\nECDSA签名: {signature}")
print(f"签名长度: {len(signature)} 字符 = {len(signature)//2} 字节")

# 4. 验证
is_valid = ecdsa_verify(message, signature, public_key)
print(f"\n签名验证: {'通过' if is_valid else '失败'}")

# 5. 使用错误公钥验证
wrong_private_key = SigningKey.generate(curve=NIST256p)
wrong_public_key = wrong_private_key.get_verifying_key()
is_valid_wrong = ecdsa_verify(message, signature, wrong_public_key)
print(f"错误公钥验证: {'通过' if is_valid_wrong else '失败'}")
\end{lstlisting}

\subsection{SM2数字签名实例（国密）}

\begin{lstlisting}[language=Python]
from gmssl import sm2

def sm2_sign(message, private_key, public_key):
    """
    使用SM2私钥签名
    """
    # 创建SM2对象
    sm2_crypt = sm2.CryptSM2(public_key=public_key, private_key=private_key)
    
    # 确保消息是字节类型
    message_bytes = message.encode('utf-8') if isinstance(message, str) else message
    
    # 签名
    signature = sm2_crypt.sign(message_bytes, private_key)
    
    return signature

def sm2_verify(message, signature, public_key):
    """
    使用SM2公钥验证签名
    """
    # 创建SM2对象
    sm2_crypt = sm2.CryptSM2(public_key=public_key, private_key="")
    
    # 确保消息是字节类型
    message_bytes = message.encode('utf-8') if isinstance(message, str) else message
    
    # 验证
    return sm2_crypt.verify(signature, message_bytes)

# ========== 完整示例 ==========

# 密钥对
private_key = '59e2e5de1614f0f830655f4b16b5903f247b701894c1f21f6f6e0631f42601e2'
public_key = '04' + '8d1d8c5df8e4286e0c8d25b602b5503e0ad4c8b6c374b8d9b8f2e8e5e5e5e5e' + \
             '9f2e5de1614f0f830655f4b16b5903f247b701894c1f21f6f6e0631f42601e2'

print("=== SM2密钥对 ===")
print(f"私钥: {private_key}")
print(f"公钥: {public_key[:20]}...{public_key[-20:]}")

# 要签名的消息
message = "这是一份使用国密SM2签名的合同"
print(f"\n原始消息: {message}")

# 签名
signature = sm2_sign(message, private_key, public_key)
print(f"\nSM2签名: {signature}")

# 验证
is_valid = sm2_verify(message, signature, public_key)
print(f"\n签名验证: {'通过' if is_valid else '失败'}")

# 篡改消息
tampered = "这是一份使用国密SM2签名的合同（已修改）"
is_valid_tampered = sm2_verify(tampered, signature, public_key)
print(f"篡改后验证: {'通过' if is_valid_tampered else '失败'}")
\end{lstlisting}

% =================== 第6章 综合应用 ===================
\section{综合应用：混合加密系统}

\subsection{为什么需要混合加密？}

\begin{keybox}[问题]
\begin{itemize}
\item 对称加密（AES/SM4）速度快，但密钥分发困难
\item 非对称加密（RSA/SM2）密钥分发容易，但速度慢且长度受限
\item \textbf{解决方案}：结合两者优点！
\end{itemize}
\end{keybox}

\subsection{混合加密方案}

\begin{lstlisting}[language=Python]
"""
混合加密方案：
1. 生成随机的对称密钥（会话密钥）
2. 用对称加密（AES）加密大数据
3. 用非对称加密（RSA）加密对称密钥
4. 发送：加密的数据 + 加密的密钥
"""

from Crypto.Cipher import AES, PKCS1_OAEP
from Crypto.PublicKey import RSA
from Crypto.Random import get_random_bytes
from Crypto.Util.Padding import pad, unpad
import base64

def hybrid_encrypt(plaintext, public_key_pem):
    """
    混合加密
    """
    # 1. 生成随机的AES密钥（会话密钥）
    session_key = get_random_bytes(32)  # 256位
    
    # 2. 用AES加密数据
    iv = get_random_bytes(16)
    cipher_aes = AES.new(session_key, AES.MODE_CBC, iv)
    
    if isinstance(plaintext, str):
        plaintext = plaintext.encode('utf-8')
    
    ciphertext = cipher_aes.encrypt(pad(plaintext, AES.block_size))
    
    # 3. 用RSA加密会话密钥
    public_key = RSA.import_key(public_key_pem)
    cipher_rsa = PKCS1_OAEP.new(public_key)
    encrypted_session_key = cipher_rsa.encrypt(session_key)
    
    # 4. 组合结果
    result = {
        'encrypted_key': base64.b64encode(encrypted_session_key).decode('utf-8'),
        'iv': base64.b64encode(iv).decode('utf-8'),
        'ciphertext': base64.b64encode(ciphertext).decode('utf-8')
    }
    
    return result

def hybrid_decrypt(encrypted_package, private_key_pem):
    """
    混合解密
    """
    # 1. 用RSA解密会话密钥
    private_key = RSA.import_key(private_key_pem)
    cipher_rsa = PKCS1_OAEP.new(private_key)
    
    encrypted_session_key = base64.b64decode(encrypted_package['encrypted_key'])
    session_key = cipher_rsa.decrypt(encrypted_session_key)
    
    # 2. 用AES解密数据
    iv = base64.b64decode(encrypted_package['iv'])
    ciphertext = base64.b64decode(encrypted_package['ciphertext'])
    
    cipher_aes = AES.new(session_key, AES.MODE_CBC, iv)
    plaintext = unpad(cipher_aes.decrypt(ciphertext), AES.block_size)
    
    return plaintext.decode('utf-8')

# ========== 完整示例 ==========

# 生成RSA密钥对
key = RSA.generate(2048)
private_key = key.export_key()
public_key = key.publickey().export_key()

# 要加密的长消息（超过RSA直接加密限制）
long_message = """
这是一段很长的消息，可能是一篇完整的文章、一份合同、
或者任何超过RSA加密长度限制的数据。RSA直接加密的长度
限制是密钥长度减去填充长度（约214字节对于2048位密钥）。

通过混合加密，我们可以：
1. 使用AES加密任意长度的数据
2. 使用RSA加密AES的密钥
3. 结合两者的优点：速度和安全性
"""

print(f"原始消息长度: {len(long_message)} 字节")

# 加密
encrypted = hybrid_encrypt(long_message, public_key)
print(f"\n加密后的会话密钥: {encrypted['encrypted_key'][:50]}...")
print(f"IV: {encrypted['iv']}")
print(f"密文长度: {len(encrypted['ciphertext'])} 字符")

# 解密
decrypted = hybrid_decrypt(encrypted, private_key)
print(f"\n解密成功: {'是' if decrypted == long_message else '否'}")
print(f"解密后长度: {len(decrypted)} 字节")
\end{lstlisting}

\begin{notebox}[运行结果]
\begin{verbatim}
原始消息长度: 312 字节

加密后的会话密钥: gAAAAABf...（Base64编码）
IV: aGVsbG8gd29ybGQ...
密文长度: 448 字符

解密成功: 是
解密后长度: 312 字节
\end{verbatim}
\end{notebox}

% =================== 第7章 总结 ===================
\section{总结}

\subsection{算法选择指南}

\begin{center}
\begin{tabular}{|l|l|l|}
\hline
\textbf{场景} & \textbf{推荐算法} & \textbf{说明} \\
\hline
大数据加密 & AES-256-GCM / SM4-CBC & 速度快，适合文件/数据 \\
\hline
密钥交换 & RSA-2048 / SM2 / ECDH & 安全分发对称密钥 \\
\hline
数据完整性 & SHA-256 / SM3 & 验证数据未被篡改 \\
\hline
数字签名 & ECDSA / SM2 / RSA-PSS & 证明身份和不可否认 \\
\hline
密码存储 & bcrypt / Argon2 / PBKDF2 & 防止彩虹表攻击 \\
\hline
\end{tabular}
\end{center}

\subsection{关键概念回顾}

\begin{enumerate}[leftmargin=1.5em]
\item \textbf{对称加密}：加密解密用同一个密钥，速度快，适合大数据
\item \textbf{非对称加密}：加密解密用不同密钥，适合密钥交换和签名
\item \textbf{Hash函数}：单向函数，用于完整性校验和签名
\item \textbf{数字签名}：证明消息来源和完整性，具有不可否认性
\item \textbf{混合加密}：结合对称和非对称加密的优点
\end{enumerate}

\subsection{安全建议}

\begin{keybox}[安全最佳实践]
\begin{enumerate}
\item 永远不要自己实现加密算法，使用经过验证的库
\item 密钥要足够长（AES-256、RSA-2048+、ECC-256+）
\item 使用安全的随机数生成器生成密钥
\item 定期更换密钥
\item 使用经过认证的加密模式（如GCM、CCM）
\item 正确处理填充和IV
\item 不要硬编码密钥在代码中
\end{enumerate}
\end{keybox}

\end{document}
